\documentclass[10pt, a4paper]{article}
\usepackage{preamble}

\title{Complex Analysis II \\
    \large Prereading Notes}
\author{Luke Phillips}
\date{January 2026}

\begin{document}

\maketitle

\newpage

\tableofcontents

\newpage

\section{The Complex Plane and Complex Functions}

\subsection{Exponential and trigonometric functions}

\begin{definition}[Complex exponential]
    We define the complex exponential function $\exp : \C \to \C$ by
    \[
    \exp(z) \coloneqq e ^ x(\cos{y} + i\sin{y}).\qquad (z = x + iy)
    \]
    short hand $\exp(z) = e ^ z$.
\end{definition}

\begin{proposition}
    We have the following properties of the complex exponential function:
    
    \begin{enumerate}[label = (\roman*)]
        \item $e ^ z \neq 0$ for all $z \in \C$.

        \item $e ^ {z_1 + z_2} = e ^ {z_1}e ^ {z_2}$.

        \item $e ^ z = 1$ if and only if $z = 2\pi ik$ for some $k \in \Z$.

        \item $e ^ {-z} = 1 / e ^ z$.

        \item $|e ^ z| = e ^ {\Re(z)}$.
    \end{enumerate}

    \begin{proof}
        Start by seeing that $e ^ {z} = e ^ xe ^ {iy}$ so
        \[
        |e ^ z| = e ^ x = e ^ {\Re(z)},\qquad \Arg(e ^ {z}) = y\mod{2\pi}.
        \]
        Since $|e ^ z| = e ^ {\Re(z)}$ we see $|e ^ z| \neq 0$ for all $z \in \C$ and $e ^ z \neq 0$ for all $z \in \C$.

        By definition
        \[
        e ^ {z_1}e ^ {z_2} = e ^ {x_1}e ^ {iy_1}e ^ {x_2}e ^ {iy_2} = e ^ {x_1}e ^ {x_2}e ^ {iy_1}e ^ {iy_2} = e ^ {x_1 + x_2}e ^ {i(y_1 + y_2)} = e ^ {z_1 + z_2}.
        \]
        This implies that
        \[
        e ^ {z}e ^ {-z} = e ^ {z - z} = e ^ 0 = 1
        \]
        and so $e ^ {-z} = 1 / e ^ z$.

        Finally,
        $e ^ z = 1 = e ^ 0e ^ {i0}$ so
        \[
        \Re(z) = 0,\qquad \Im(z) = 0\mod{2\pi}
        \]
        in other words $z = 2\pi ik$ for some $k \in \Z$.
    \end{proof}
\end{proposition}

\begin{corollary}
    We have $\exp(2\pi i) = 1$ and $\exp(\pi i) = -1$.
    The latter being Euler's formula.
\end{corollary}

\begin{corollary}
    The complex exponential function is $2\pi i$-periodic,
    $\exp(z + 2k\pi i) = \exp(z)$ for any $k \in \Z$.
\end{corollary}




\end{document}